\section{给个人和组织的行动建议}

\subsection{个人层：把 AI 当作长期能力放大器}

访谈后段最具行动性的观点是这句反复出现的话：\textit{``How can I use AI to do my job better?''} Jensen 认为这应成为跨职业、跨年龄的共同问题。

\begin{knowledgebox}{可执行的三步法}
\begin{enumerate}
    \item 每天固定一个高频任务，用 AI 做 first draft（文档、分析、代码、检索）。
    \item 用专家判断做 second pass（纠错、约束、风格、责任边界）。
    \item 记录“提示词与结果”的可复用模板，形成个人 workflow 资产。
\end{enumerate}
\end{knowledgebox}

\begin{importantbox}{学习门槛变化}
Jensen 的观察很实用：对于从未接触电脑的人，传统软件学习门槛很高；而 LLM 可以直接通过自然语言指导新手完成入门。智能系统降低了使用智能系统的门槛。
\end{importantbox}

\subsection{组织层：从“买工具”转向“重构流程”}

组织采用 AI 失败的常见原因是只采购工具，不重构流程。可行路径应是：
\begin{itemize}
    \item 选 2-3 条高价值流程做闭环试点。
    \item 明确人机分工与审批责任。
    \item 建立质量、时延、成本三维指标，滚动迭代。
\end{itemize}

\begin{warningbox}{避免“演示幻觉”}
许多团队停留在 demo 成功阶段，却没有进入生产级可靠性验证。真正的价值来自持续可重复的流程优化，而不是一次性炫技展示。
\end{warningbox}

\subsection{本章小结}
个人应把 AI 视作“能力复利工具”，组织应把 AI 视作“流程重构工程”。这比单纯追逐模型参数更可持续。

